\documentclass[11pt]{article}
% =============================================================================
%  Shared article preamble — tutorials, homework, solutions.
%  Same fonts, palette and notation as the Beamer decks (see preamble.tex).
%      \documentclass[11pt]{article}
%      % =============================================================================
%  Shared article preamble — tutorials, homework, solutions.
%  Same fonts, palette and notation as the Beamer decks (see preamble.tex).
%      \documentclass[11pt]{article}
%      % =============================================================================
%  Shared article preamble — tutorials, homework, solutions.
%  Same fonts, palette and notation as the Beamer decks (see preamble.tex).
%      \documentclass[11pt]{article}
%      \input{../../../style/preamble_article}
% =============================================================================
\usepackage[margin=1in]{geometry}
\usepackage{helvet}
\usepackage[default]{lato}
\usepackage{mathpazo}
\usepackage[utf8]{inputenc}
\usepackage{amsmath,amssymb,amsthm}
\usepackage{mathtools}
\usepackage{graphicx}
\usepackage{booktabs}
\usepackage{multirow}
\usepackage{enumitem}
\usepackage{xcolor}
\usepackage{tcolorbox}
\tcbuselibrary{skins,breakable}
\usepackage{listings}
\usepackage[ruled,vlined]{algorithm2e}
\usepackage{hyperref}
\usepackage{fancyhdr}
\usepackage{titlesec}

\definecolor{blue}{RGB}{0,114,178}
\definecolor{red}{RGB}{213,94,0}
\definecolor{yellow}{RGB}{240,228,66}
\definecolor{green}{RGB}{0,158,115}
\hypersetup{colorlinks=true, linkcolor=blue, urlcolor=blue, citecolor=blue}

\titleformat{\section}{\Large\bfseries\color{blue}}{\thesection}{0.6em}{}
\titleformat{\subsection}{\large\bfseries}{\thesubsection}{0.6em}{}
\pagestyle{fancy}\fancyhf{}
\rhead{\footnotesize ENEI -- INEI \,\textperiodcentered\, PEU-CD 2026}
\lfoot{\footnotesize Machine Learning I \& II \,\textperiodcentered\, Alexander Quispe}
\rfoot{\footnotesize\thepage}
\renewcommand{\headrulewidth}{0.3pt}

\newtcolorbox{derivation}[1]{enhanced, breakable, colback=blue!3, colframe=blue, fonttitle=\bfseries,
  title={#1}, boxrule=0.6pt, left=6pt, right=6pt, top=4pt, bottom=4pt, arc=2pt}
\newtcolorbox{claim}[1]{enhanced, breakable, colback=white, colframe=black, fonttitle=\bfseries,
  title={#1}, boxrule=0.6pt, left=6pt, right=6pt, top=4pt, bottom=4pt, arc=2pt}
\newtcolorbox{keypoint}{enhanced, breakable, colback=green!6, colframe=green, boxrule=0.6pt,
  left=6pt, right=6pt, top=4pt, bottom=4pt, arc=2pt}
\newtcolorbox{caution}{enhanced, breakable, colback=red!5, colframe=red, boxrule=0.6pt,
  left=6pt, right=6pt, top=4pt, bottom=4pt, arc=2pt}
\newtcolorbox{task}[1]{enhanced, breakable, colback=yellow!12, colframe=yellow!60!black, fonttitle=\bfseries,
  title={#1}, boxrule=0.6pt, left=6pt, right=6pt, top=4pt, bottom=4pt, arc=2pt}

\lstset{
  basicstyle=\ttfamily\small, keywordstyle=\color{blue}\bfseries,
  commentstyle=\color{green}, stringstyle=\color{red},
  showstringspaces=false, breaklines=true, frame=single, rulecolor=\color{black!30},
  numbers=none, xleftmargin=6pt, framexleftmargin=6pt, language=Python
}
\newtheorem{theorem}{Theorem}
\newtheorem{lemma}{Lemma}
\newtheorem{proposition}{Proposition}
\theoremstyle{definition}
\newtheorem{definition}{Definition}
\newtheorem{exercise}{Exercise}

\input{../../../style/macros}

% =============================================================================
\usepackage[margin=1in]{geometry}
\usepackage{helvet}
\usepackage[default]{lato}
\usepackage{mathpazo}
\usepackage[utf8]{inputenc}
\usepackage{amsmath,amssymb,amsthm}
\usepackage{mathtools}
\usepackage{graphicx}
\usepackage{booktabs}
\usepackage{multirow}
\usepackage{enumitem}
\usepackage{xcolor}
\usepackage{tcolorbox}
\tcbuselibrary{skins,breakable}
\usepackage{listings}
\usepackage[ruled,vlined]{algorithm2e}
\usepackage{hyperref}
\usepackage{fancyhdr}
\usepackage{titlesec}

\definecolor{blue}{RGB}{0,114,178}
\definecolor{red}{RGB}{213,94,0}
\definecolor{yellow}{RGB}{240,228,66}
\definecolor{green}{RGB}{0,158,115}
\hypersetup{colorlinks=true, linkcolor=blue, urlcolor=blue, citecolor=blue}

\titleformat{\section}{\Large\bfseries\color{blue}}{\thesection}{0.6em}{}
\titleformat{\subsection}{\large\bfseries}{\thesubsection}{0.6em}{}
\pagestyle{fancy}\fancyhf{}
\rhead{\footnotesize ENEI -- INEI \,\textperiodcentered\, PEU-CD 2026}
\lfoot{\footnotesize Machine Learning I \& II \,\textperiodcentered\, Alexander Quispe}
\rfoot{\footnotesize\thepage}
\renewcommand{\headrulewidth}{0.3pt}

\newtcolorbox{derivation}[1]{enhanced, breakable, colback=blue!3, colframe=blue, fonttitle=\bfseries,
  title={#1}, boxrule=0.6pt, left=6pt, right=6pt, top=4pt, bottom=4pt, arc=2pt}
\newtcolorbox{claim}[1]{enhanced, breakable, colback=white, colframe=black, fonttitle=\bfseries,
  title={#1}, boxrule=0.6pt, left=6pt, right=6pt, top=4pt, bottom=4pt, arc=2pt}
\newtcolorbox{keypoint}{enhanced, breakable, colback=green!6, colframe=green, boxrule=0.6pt,
  left=6pt, right=6pt, top=4pt, bottom=4pt, arc=2pt}
\newtcolorbox{caution}{enhanced, breakable, colback=red!5, colframe=red, boxrule=0.6pt,
  left=6pt, right=6pt, top=4pt, bottom=4pt, arc=2pt}
\newtcolorbox{task}[1]{enhanced, breakable, colback=yellow!12, colframe=yellow!60!black, fonttitle=\bfseries,
  title={#1}, boxrule=0.6pt, left=6pt, right=6pt, top=4pt, bottom=4pt, arc=2pt}

\lstset{
  basicstyle=\ttfamily\small, keywordstyle=\color{blue}\bfseries,
  commentstyle=\color{green}, stringstyle=\color{red},
  showstringspaces=false, breaklines=true, frame=single, rulecolor=\color{black!30},
  numbers=none, xleftmargin=6pt, framexleftmargin=6pt, language=Python
}
\newtheorem{theorem}{Theorem}
\newtheorem{lemma}{Lemma}
\newtheorem{proposition}{Proposition}
\theoremstyle{definition}
\newtheorem{definition}{Definition}
\newtheorem{exercise}{Exercise}

% =============================================================================
%  Notation — one convention for all lectures
%
%  scalars      x, y, w           plain italic
%  vectors      \vx \vy \vw       bold lower-case
%  matrices     \vX \vW \vSigma   bold upper-case
%  parameters   \vbeta \vtheta    bold Greek
%  transpose    \vX\T             \top
% =============================================================================
\newcommand{\R}{\mathbb{R}}
\newcommand{\E}{\mathbb{E}}
\newcommand{\Prob}{\mathbb{P}}
\DeclareMathOperator{\Var}{Var}
\DeclareMathOperator{\Cov}{Cov}
\DeclareMathOperator{\tr}{tr}
\DeclareMathOperator{\diag}{diag}
\DeclareMathOperator{\rank}{rank}
\DeclareMathOperator{\sign}{sign}
\DeclareMathOperator{\col}{col}
\DeclareMathOperator{\RSS}{RSS}
\DeclareMathOperator*{\argmin}{arg\,min}
\DeclareMathOperator*{\argmax}{arg\,max}
\newcommand{\norm}[1]{\left\lVert #1 \right\rVert}
\newcommand{\abs}[1]{\left\lvert #1 \right\rvert}
\newcommand{\inner}[2]{\left\langle #1,\, #2 \right\rangle}
\newcommand{\ind}[1]{\mathbf{1}\!\left\{#1\right\}}
\newcommand{\T}{^{\top}}
\newcommand{\inv}{^{-1}}
\newcommand{\half}{\tfrac{1}{2}}
\newcommand{\eps}{\varepsilon}
\newcommand{\Dset}{\mathcal{D}}
\newcommand{\Hset}{\mathcal{H}}
\newcommand{\Lcal}{\mathcal{L}}
\newcommand{\Ncal}{\mathcal{N}}

% bold vectors
\newcommand{\vx}{\mathbf{x}}   \newcommand{\vy}{\mathbf{y}}   \newcommand{\vz}{\mathbf{z}}
\newcommand{\vw}{\mathbf{w}}   \newcommand{\vu}{\mathbf{u}}   \newcommand{\vv}{\mathbf{v}}
\newcommand{\vb}{\mathbf{b}}   \newcommand{\vh}{\mathbf{h}}   \newcommand{\vr}{\mathbf{r}}
\newcommand{\vc}{\mathbf{c}}   \newcommand{\vf}{\mathbf{f}}   \newcommand{\vg}{\mathbf{g}}
\newcommand{\vp}{\mathbf{p}}   \newcommand{\vq}{\mathbf{q}}   \newcommand{\vs}{\mathbf{s}}
\newcommand{\va}{\mathbf{a}}   \newcommand{\vd}{\mathbf{d}}   \newcommand{\vt}{\mathbf{t}}
\newcommand{\vk}{\mathbf{k}}   \newcommand{\vm}{\mathbf{m}}   \newcommand{\vn}{\mathbf{n}}
\newcommand{\vzero}{\mathbf{0}} \newcommand{\vone}{\mathbf{1}}
% bold matrices
\newcommand{\vX}{\mathbf{X}}   \newcommand{\vY}{\mathbf{Y}}   \newcommand{\vZ}{\mathbf{Z}}
\newcommand{\vW}{\mathbf{W}}   \newcommand{\vH}{\mathbf{H}}   \newcommand{\vI}{\mathbf{I}}
\newcommand{\vA}{\mathbf{A}}   \newcommand{\vB}{\mathbf{B}}   \newcommand{\vC}{\mathbf{C}}
\newcommand{\vD}{\mathbf{D}}   \newcommand{\vP}{\mathbf{P}}   \newcommand{\vQ}{\mathbf{Q}}
\newcommand{\vU}{\mathbf{U}}   \newcommand{\vV}{\mathbf{V}}   \newcommand{\vS}{\mathbf{S}}
\newcommand{\vM}{\mathbf{M}}   \newcommand{\vK}{\mathbf{K}}   \newcommand{\vG}{\mathbf{G}}
% bold Greek
\newcommand{\vbeta}{\boldsymbol{\beta}}     \newcommand{\valpha}{\boldsymbol{\alpha}}
\newcommand{\vtheta}{\boldsymbol{\theta}}   \newcommand{\vmu}{\boldsymbol{\mu}}
\newcommand{\vSigma}{\boldsymbol{\Sigma}}   \newcommand{\vLambda}{\boldsymbol{\Lambda}}
\newcommand{\veps}{\boldsymbol{\varepsilon}} \newcommand{\vphi}{\boldsymbol{\phi}}
\newcommand{\vgamma}{\boldsymbol{\gamma}}   \newcommand{\vlambda}{\boldsymbol{\lambda}}
\newcommand{\vdelta}{\boldsymbol{\delta}}   \newcommand{\vnu}{\boldsymbol{\nu}}
\newcommand{\vxi}{\boldsymbol{\xi}}         \newcommand{\vpi}{\boldsymbol{\pi}}
\newcommand{\vomega}{\boldsymbol{\omega}}



% =============================================================================
\usepackage[margin=1in]{geometry}
\usepackage{helvet}
\usepackage[default]{lato}
\usepackage{mathpazo}
\usepackage[utf8]{inputenc}
\usepackage{amsmath,amssymb,amsthm}
\usepackage{mathtools}
\usepackage{graphicx}
\usepackage{booktabs}
\usepackage{multirow}
\usepackage{enumitem}
\usepackage{xcolor}
\usepackage{tcolorbox}
\tcbuselibrary{skins,breakable}
\usepackage{listings}
\usepackage[ruled,vlined]{algorithm2e}
\usepackage{hyperref}
\usepackage{fancyhdr}
\usepackage{titlesec}

\definecolor{blue}{RGB}{0,114,178}
\definecolor{red}{RGB}{213,94,0}
\definecolor{yellow}{RGB}{240,228,66}
\definecolor{green}{RGB}{0,158,115}
\hypersetup{colorlinks=true, linkcolor=blue, urlcolor=blue, citecolor=blue}

\titleformat{\section}{\Large\bfseries\color{blue}}{\thesection}{0.6em}{}
\titleformat{\subsection}{\large\bfseries}{\thesubsection}{0.6em}{}
\pagestyle{fancy}\fancyhf{}
\rhead{\footnotesize ENEI -- INEI \,\textperiodcentered\, PEU-CD 2026}
\lfoot{\footnotesize Machine Learning I \& II \,\textperiodcentered\, Alexander Quispe}
\rfoot{\footnotesize\thepage}
\renewcommand{\headrulewidth}{0.3pt}

\newtcolorbox{derivation}[1]{enhanced, breakable, colback=blue!3, colframe=blue, fonttitle=\bfseries,
  title={#1}, boxrule=0.6pt, left=6pt, right=6pt, top=4pt, bottom=4pt, arc=2pt}
\newtcolorbox{claim}[1]{enhanced, breakable, colback=white, colframe=black, fonttitle=\bfseries,
  title={#1}, boxrule=0.6pt, left=6pt, right=6pt, top=4pt, bottom=4pt, arc=2pt}
\newtcolorbox{keypoint}{enhanced, breakable, colback=green!6, colframe=green, boxrule=0.6pt,
  left=6pt, right=6pt, top=4pt, bottom=4pt, arc=2pt}
\newtcolorbox{caution}{enhanced, breakable, colback=red!5, colframe=red, boxrule=0.6pt,
  left=6pt, right=6pt, top=4pt, bottom=4pt, arc=2pt}
\newtcolorbox{task}[1]{enhanced, breakable, colback=yellow!12, colframe=yellow!60!black, fonttitle=\bfseries,
  title={#1}, boxrule=0.6pt, left=6pt, right=6pt, top=4pt, bottom=4pt, arc=2pt}

\lstset{
  basicstyle=\ttfamily\small, keywordstyle=\color{blue}\bfseries,
  commentstyle=\color{green}, stringstyle=\color{red},
  showstringspaces=false, breaklines=true, frame=single, rulecolor=\color{black!30},
  numbers=none, xleftmargin=6pt, framexleftmargin=6pt, language=Python
}
\newtheorem{theorem}{Theorem}
\newtheorem{lemma}{Lemma}
\newtheorem{proposition}{Proposition}
\theoremstyle{definition}
\newtheorem{definition}{Definition}
\newtheorem{exercise}{Exercise}

% =============================================================================
%  Notation — one convention for all lectures
%
%  scalars      x, y, w           plain italic
%  vectors      \vx \vy \vw       bold lower-case
%  matrices     \vX \vW \vSigma   bold upper-case
%  parameters   \vbeta \vtheta    bold Greek
%  transpose    \vX\T             \top
% =============================================================================
\newcommand{\R}{\mathbb{R}}
\newcommand{\E}{\mathbb{E}}
\newcommand{\Prob}{\mathbb{P}}
\DeclareMathOperator{\Var}{Var}
\DeclareMathOperator{\Cov}{Cov}
\DeclareMathOperator{\tr}{tr}
\DeclareMathOperator{\diag}{diag}
\DeclareMathOperator{\rank}{rank}
\DeclareMathOperator{\sign}{sign}
\DeclareMathOperator{\col}{col}
\DeclareMathOperator{\RSS}{RSS}
\DeclareMathOperator*{\argmin}{arg\,min}
\DeclareMathOperator*{\argmax}{arg\,max}
\newcommand{\norm}[1]{\left\lVert #1 \right\rVert}
\newcommand{\abs}[1]{\left\lvert #1 \right\rvert}
\newcommand{\inner}[2]{\left\langle #1,\, #2 \right\rangle}
\newcommand{\ind}[1]{\mathbf{1}\!\left\{#1\right\}}
\newcommand{\T}{^{\top}}
\newcommand{\inv}{^{-1}}
\newcommand{\half}{\tfrac{1}{2}}
\newcommand{\eps}{\varepsilon}
\newcommand{\Dset}{\mathcal{D}}
\newcommand{\Hset}{\mathcal{H}}
\newcommand{\Lcal}{\mathcal{L}}
\newcommand{\Ncal}{\mathcal{N}}

% bold vectors
\newcommand{\vx}{\mathbf{x}}   \newcommand{\vy}{\mathbf{y}}   \newcommand{\vz}{\mathbf{z}}
\newcommand{\vw}{\mathbf{w}}   \newcommand{\vu}{\mathbf{u}}   \newcommand{\vv}{\mathbf{v}}
\newcommand{\vb}{\mathbf{b}}   \newcommand{\vh}{\mathbf{h}}   \newcommand{\vr}{\mathbf{r}}
\newcommand{\vc}{\mathbf{c}}   \newcommand{\vf}{\mathbf{f}}   \newcommand{\vg}{\mathbf{g}}
\newcommand{\vp}{\mathbf{p}}   \newcommand{\vq}{\mathbf{q}}   \newcommand{\vs}{\mathbf{s}}
\newcommand{\va}{\mathbf{a}}   \newcommand{\vd}{\mathbf{d}}   \newcommand{\vt}{\mathbf{t}}
\newcommand{\vk}{\mathbf{k}}   \newcommand{\vm}{\mathbf{m}}   \newcommand{\vn}{\mathbf{n}}
\newcommand{\vzero}{\mathbf{0}} \newcommand{\vone}{\mathbf{1}}
% bold matrices
\newcommand{\vX}{\mathbf{X}}   \newcommand{\vY}{\mathbf{Y}}   \newcommand{\vZ}{\mathbf{Z}}
\newcommand{\vW}{\mathbf{W}}   \newcommand{\vH}{\mathbf{H}}   \newcommand{\vI}{\mathbf{I}}
\newcommand{\vA}{\mathbf{A}}   \newcommand{\vB}{\mathbf{B}}   \newcommand{\vC}{\mathbf{C}}
\newcommand{\vD}{\mathbf{D}}   \newcommand{\vP}{\mathbf{P}}   \newcommand{\vQ}{\mathbf{Q}}
\newcommand{\vU}{\mathbf{U}}   \newcommand{\vV}{\mathbf{V}}   \newcommand{\vS}{\mathbf{S}}
\newcommand{\vM}{\mathbf{M}}   \newcommand{\vK}{\mathbf{K}}   \newcommand{\vG}{\mathbf{G}}
% bold Greek
\newcommand{\vbeta}{\boldsymbol{\beta}}     \newcommand{\valpha}{\boldsymbol{\alpha}}
\newcommand{\vtheta}{\boldsymbol{\theta}}   \newcommand{\vmu}{\boldsymbol{\mu}}
\newcommand{\vSigma}{\boldsymbol{\Sigma}}   \newcommand{\vLambda}{\boldsymbol{\Lambda}}
\newcommand{\veps}{\boldsymbol{\varepsilon}} \newcommand{\vphi}{\boldsymbol{\phi}}
\newcommand{\vgamma}{\boldsymbol{\gamma}}   \newcommand{\vlambda}{\boldsymbol{\lambda}}
\newcommand{\vdelta}{\boldsymbol{\delta}}   \newcommand{\vnu}{\boldsymbol{\nu}}
\newcommand{\vxi}{\boldsymbol{\xi}}         \newcommand{\vpi}{\boldsymbol{\pi}}
\newcommand{\vomega}{\boldsymbol{\omega}}




\title{\textcolor{blue}{Lab 4 --- Clustering, PCA and the ML I Project}\\[4pt]
\large Machine Learning I \,\textperiodcentered\, Tutorial}
\author{Rodrigo Grijalba (teaching assistant) \,\textperiodcentered\, Alexander Quispe (instructor)}
\date{Wednesday, September 30, 2026 \,\textperiodcentered\, 19:00--22:00}

\begin{document}
\maketitle

\begin{keypoint}
\textbf{Goal of the session.} Implement $k$-means and watch its objective fall exactly as Lecture 6's
proof says it must; verify Ward's merge formula on real clusters; compute principal components two ways
and confirm they agree; check the variance--reconstruction identity numerically; and close ML I with
the project.
\end{keypoint}

\section*{Materials}
\begin{itemize}[nosep]
  \item \texttt{lab\_04\_student.ipynb} / \texttt{lab\_04\_solution.ipynb}.
  \item Sources: CS 155 Problem Set 5 (PCA, SVD, low-rank approximation).
  \item Data: \texttt{make\_blobs}, \texttt{load\_digits} ($1797\times64$), \texttt{load\_diabetes} --- all bundled.
\end{itemize}

\section{$k$-Means (40 min)}

\begin{task}{Task 1 --- Lloyd's algorithm, and the lecture example}
Write \texttt{kmeans(X, mu0, iters)} that alternates the two half-steps of Lecture 6 and records $J$ after
each. Run it on the six points $A(1,1)$, $B(1.5,2)$, $C(3,3.5)$, $D(5,7)$, $E(3.5,5)$, $F(4.5,5)$ from
$\vmu_1=A$, $\vmu_2=D$. You must get $J=22$ after the first assignment and $J=9.167$ at convergence, with
final centroids $(1.833,2.167)$ and $(4.333,5.667)$. Assert that the recorded $J$ sequence is non-increasing.
\end{task}

\begin{task}{Task 2 --- local minima, and $k$-means++}
On \texttt{make\_blobs(300, centers=4, cluster\_std=0.9)}, run your \texttt{kmeans} from 50 random
initializations (centroids drawn from the data) and histogram the final $J$. How many distinct local minima
do you find? Then implement $k$-means++ seeding (the $D^2$ sampling rule) and repeat. Compare the two histograms
and the best $J$ from each. Finally confirm \texttt{sklearn.cluster.KMeans} with your best seeds reaches the same $J$.
\end{task}

\begin{task}{Task 3 --- choosing $K$}
For $K=1,\dots,10$ on the same data, plot $J$ (elbow) and the mean silhouette
(\texttt{sklearn.metrics.silhouette\_score}). Do they agree on $K=4$? Then add one far outlier to the data
and repeat. What happens to the elbow, and why (Lecture 6, ``What $k$-Means Assumes'')?
\end{task}

\section{Hierarchical Clustering (25 min)}

\begin{task}{Task 4 --- Ward's formula, verified}
Take two clusters $A$ and $B$ from Task 2's data (any two of the four). Compute the within-cluster sum of
squares $J$ with them separate and with them merged, and check that the increase equals
$\frac{n_An_B}{n_A+n_B}\norm{\vmu_A-\vmu_B}^2$ to floating-point precision. Then compute
\texttt{scipy.cluster.hierarchy.linkage(X, "ward")} and confirm its last merge height, squared, equals the
$\Delta J$ of merging the final two clusters (scipy reports $\sqrt{2\Delta J}$; check its documentation and
reconcile the factor).
\end{task}

\begin{task}{Task 5 --- four linkages, one dendrogram each}
Draw the dendrograms for single, complete, average and Ward linkage on 40 points from three blobs.
Cut each at three clusters and compare the partitions to the truth with the adjusted Rand index. Which linkage
fails, and on what kind of data would it be the right choice?
\end{task}

\section{PCA (50 min)}

\begin{task}{Task 6 --- two computations, one answer}
Centre the digits data $\vX$ ($1797\times64$).
\begin{enumerate}[nosep]
  \item Eigendecompose $\vS=\vX\T\vX/N$ with \texttt{np.linalg.eigh}; sort descending.
  \item Compute the thin SVD of $\vX$; verify $\vV$ matches the eigenvectors (up to sign) and $\lambda_j=d_j^2/N$.
  \item Verify the Ky Fan claim: for $k=5$, $\tr(\vV_k\T\vS\vV_k)=\lambda_1+\dots+\lambda_5$, and it exceeds
        $\tr(\vW\T\vS\vW)$ for 100 random orthonormal $\vW\in\R^{64\times5}$ (orthonormalize with \texttt{np.linalg.qr}).
\end{enumerate}
\end{task}

\begin{task}{Task 7 --- variance $=$ total $-$ reconstruction}
For $k=1,2,4,8,16,32,64$ compute (i) the mean squared reconstruction error $\frac1N\sum_i\norm{\vx_i-\vV_k\vV_k\T\vx_i}^2$,
(ii) $\tr\vS-\sum_{j\le k}\lambda_j$. They must agree. Plot the proportion of variance explained against $k$;
report the smallest $k$ reaching $90\%$.
\end{task}

\begin{task}{Task 8 --- eigendigits}
Reshape the first eight principal components to $8\times8$ and display them. Reconstruct one digit with
$k=1,2,4,8,16,32,64$ components and display the sequence. At which $k$ can you read the digit?
\end{task}

\begin{task}{Task 9 --- hard vs.\ soft: PCR against ridge}
On the diabetes data (standardized), compare by 10-fold CV:
principal-components regression with $k=1,\dots,10$ components, and ridge with $\lambda$ on a log grid.
Plot both CV curves against effective degrees of freedom ($k$ for PCR; $\sum_jd_j^2/(d_j^2+\lambda)$ for ridge).
Lecture 6 called PCR a hard threshold and ridge a soft one; which wins here, and by how much?
\end{task}

\section{The ML I Project (due today, 23:59)}

One dataset, the whole pipeline, a short report. Choose \emph{one}:
\begin{itemize}[nosep]
  \item \texttt{load\_breast\_cancer} --- binary classification, $569\times30$.
  \item \texttt{load\_digits} --- 10-class classification, $1797\times64$.
  \item \texttt{load\_diabetes} --- regression, $442\times10$.
  \item The Messidor retinopathy set from CS 155 Set 3, at\\ \texttt{reference/caltech\_cs155/problem\_sets/set3/messidor\_features.arff}.
\end{itemize}

\begin{task}{Deliverable}
A notebook and a PDF report of at most 4 pages, containing:
\begin{enumerate}[nosep]
  \item \textbf{Setup} --- the task, the loss you are optimizing, the metric you will report, and why they differ if they do.
        The test split, sealed before anything else.
  \item \textbf{Baseline} --- the trivial predictor, and its number.
  \item \textbf{Three models} --- one regularized linear model (ridge/lasso/logistic), one random forest, one
        gradient-boosting model, each tuned by cross-validation with the search grid stated. Standardization inside folds.
  \item \textbf{Evaluation} --- the chosen metric on the test set for all four, once. For classification, also a
        confusion matrix and ROC/AUC; for regression, residual plots.
  \item \textbf{One derivation} --- pick any result from Lectures 1--6 that your pipeline relies on (the optimism
        formula, the bagging variance, the soft-threshold rule, \dots) and re-derive it in your own words in half a page.
  \item \textbf{What you would do next} --- three sentences.
\end{enumerate}
\end{task}

\textbf{Rubric} (out of 20): setup and protocol 4, baseline and models 6, evaluation 4, derivation 4, clarity 2.
A pipeline that touches the test set more than once loses the 4 protocol points, whatever its numbers.

\section{Exercises (optional; not graded)}
\begin{exercise}
For standardized 2-D data with correlation $\rho>0$, show that
\[\vS=\begin{pmatrix}1&\rho\\\rho&1\end{pmatrix}\]
has eigenvectors $(1,1)/\sqrt2$ and $(1,-1)/\sqrt2$ with eigenvalues $1+\rho$ and $1-\rho$. Verify numerically.
\end{exercise}
\begin{exercise}
Implement power iteration for the top eigenvector of $\vS$: $\vu\leftarrow\vS\vu/\norm{\vS\vu}$. Show that
the angle to \texttt{eigh}'s answer decays like $(\lambda_2/\lambda_1)^t$ on the digits data.
\end{exercise}

\section*{Submission}
Notebook \texttt{lab04\_lastname\_firstname.ipynb}, and the project as\\ \texttt{project\_lastname\_firstname.ipynb} $+$ \texttt{.pdf}, to the course drive before Wednesday, September 30, 23:59.

\end{document}
